% ---
% Agradecimentos

% \begin{agradecimentos}

%Agradeço, em primeiro lugar, aos meus pais, Anailton e Ana Maria, por me darem a vida e por serem parte essencial da minha história. Sem eles, eu não estaria neste mundo, nem teria chegado até este momento.

%Agradeço ao meu orientador, professor Degas, por todo apoio, correções, incentivo e acompanhamento ao longo deste trabalho. Suas contribuições foram fundamentais para o amadurecimento desta pesquisa e para que eu pudesse concluir esta etapa com mais segurança. Agradeço também à professora Ximena pelo acompanhamento, ao professor Jorge pelo incentivo e apoio, e a todos os demais professores da UESC que, de diferentes formas, tiveram papel fundamental na minha formação acadêmica e pessoal.

%Agradeço à minha avó Guiomar, que, mesmo sem ter tido acesso à leitura e sem compreender os detalhes da minha graduação, sempre me incentivou a continuar. Seu apoio simples, firme e amoroso também faz parte desta conquista.

%Agradeço aos meus irmãos, Raul, Lara, Pablo e Amora, pela presença, pelo afeto e por fazerem parte da minha caminhada.

%Agradeço aos meus amigos que estiveram comigo durante esse percurso. A Ezequiel, por ser a pessoa que mais me apoia neste mundo, meu porto seguro e uma fonte constante de felicidade. A Ariane, pela parceria e por sempre acreditar em um potencial que muitas vezes eu mesmo não conseguia enxergar. A Bruno, por sempre me apoiar, me incentivar e acreditar em mim. A Vitor, por tratar minhas dificuldades como algo que eu seria capaz de resolver, sem dúvidas.

%Agradeço também a todos os outros amigos que não foram mencionados nominalmente, aos não tão amigos assim e à vida. Cada experiência, boa ou difícil, contribuiu de alguma forma para que eu me tornasse quem sou, melhorasse e continuasse tentando melhorar.

%\end{agradecimentos}
% 
\begin{agradecimentos}
    
Agradeço, em primeiro lugar, aos meus pais, Anailton e Ana Maria, por me darem a vida e por serem parte essencial da minha história. Sem eles, eu não estaria neste mundo.

À minha avó, Guiomar, cujo amor e sabedoria transcendem qualquer barreira. Mesmo sem compreender as complexidades técnicas da minha graduação, a senhora foi uma fonte inesgotável de incentivo, ensinando-me que a maior força muitas vezes reside na simplicidade de acreditar.

Aos meus irmãos, Raul, Lara, Pablo e Amora, por fazerem parte da minha história, por dividirem a vida comigo e por serem parte fundamental da minha jornada.

Ao meu orientador, professor Degas, por todo o apoio, pela paciência nas correções e pelo incentivo constante que foram o alicerce deste trabalho. À professora Ximena, pelo acompanhamento cuidadoso durante o processo, e ao professor Jorge, pelo apoio e incentivo fundamentais. Estendo minha gratidão a todos os professores da Universidade Estadual de Santa Cruz (UESC), que tiveram um papel decisivo na minha formação profissional e pessoal ao longo destes anos.

Ao Ezequiel, meu porto seguro e motivo de tantas alegrias diárias. Você é a pessoa que mais me apoia neste mundo, e a sua presença ao meu lado tornou o peso desta graduação muito mais leve.

À Ariane, minha amiga e parceira, que sempre enxergou e acreditou em um potencial que eu mesma, em muitos momentos de incerteza, hesitei em ver.

Aos grandes parceiros de jornada, Bruno, por sempre me apoiar, me incentivar e por não me deixar desistir; e Vitor, por sempre tratar as dificuldades e os problemas técnicos como algo que eu, sem a menor dúvida, seria capaz de resolver. A confiança de vocês foi um combustível essencial.

Aos demais amigos que não foram citados nominalmente, mas que sabem da importância que têm na minha vida. E até mesmo aos "não tão amigos assim", cujas passagens cruzaram o meu caminho. Deixo, por fim, o meu agradecimento à própria vida: cada pessoa, cada desafio e cada experiência me moldaram de alguma forma, ensinando-me a ser resiliente e a buscar ser alguém melhor a cada day.

\end{agradecimentos}